% ============================== UNIT 5 =====================================
\unitheading{5}{Atomic properties and periodicity}

\textit{Handout chapter 5, sections 5.1--5.5 (pages 16--20).} Five properties of atoms (mass, radius,
ionisation energy, polarisability, electronegativity) and how they vary across the periodic table. Most
trends are explained by two ideas: the shell $n$ of the valence electrons and the effective nuclear
charge they feel.

\partheading{Concepts in detail}

\sect{Key idea}{Effective nuclear charge, the key to every trend}
\begin{concept}[Screening and effective nuclear charge $Z_\mathrm{eff}$.]
A valence electron is attracted by the $Z$ protons but repelled by the other electrons, especially the
core electrons between it and the nucleus. It feels only an effective charge
$Z_\mathrm{eff} = Z - \sigma$, where $\sigma$ (the screening constant) measures the shielding by the
other electrons.
\end{concept}
\begin{itemize}
  \item \textbf{Across a period} (left to right): one proton and one valence electron are added each
        time, but electrons of the same shell screen each other poorly. $Z_\mathrm{eff}$ increases
        (roughly $+0.65$ per step), the valence shell $n$ stays the same: valence electrons are held more
        tightly.
  \item \textbf{Down a group}: $Z_\mathrm{eff}$ changes little, but $n$ increases by one: the valence
        electrons are in a larger shell, further from the nucleus, and held more loosely.
\end{itemize}
\begin{going}
Slater's rules estimate $\sigma$ for an $n$s or $n$p electron: each other electron of the same shell
contributes 0.35 (0.30 in 1s), each electron of shell $n-1$ contributes 0.85, each electron of lower
shells 1.00. Example, Cl ($[\mathrm{Ne}]\,3\mathrm s^2 3\mathrm p^5$): $\sigma = 6\times0.35 + 8\times0.85 +
2\times1.00 = 10.9$ (the 6 other valence electrons, then the 8 electrons of shell 2, then the 2 of shell 1),
so $Z_\mathrm{eff} = 17 - 10.9 = 6.1$. Na: $\sigma = 8\times0.85 + 2 = 8.8$, $Z_\mathrm{eff} = 2.2$.
\end{going}

\sect{5.1}{Atomic mass}
\begin{concept}[Atomic mass.]
The mass of one atom, in atomic mass units (amu). It is close to the mass number $A$ (protons +
neutrons, each $\approx1$~amu); the electrons are negligible (Unit 1). Numerically, the atomic mass in amu
equals the molar mass in $\mathrm{g\,mol^{-1}}$.
\end{concept}
\textbf{Measurement}: the mass spectrometer ionises the atoms, accelerates them and separates them by
mass-to-charge ratio in a magnetic or electric field. It gives the mass of each isotope and its
abundance.
\begin{handout}{bullets of 5.1.2, trend in the periodic table}
\begin{itemize}
  \item The atomic mass increases with $Z$ along a period: each step adds a proton and usually one or
        more neutrons.
  \item It increases down a group (a whole period of nucleons is added). There are a few
        \emph{inversions}, where mass and $Z$ are not in the same order: Ar/K, Co/Ni, Te/I.
\end{itemize}
\end{handout}
\begin{concept}[Average atomic mass.]
Most elements are mixtures of isotopes. The tabulated atomic mass is the average of the isotope
masses weighted by their natural abundances $x_i$ (fractions, $\sum x_i = 1$):
$\boxed{\bar m = \sum_i x_i\,m_i}$.
\end{concept}
\begin{handout}{chlorine, argon and potassium example}
\begin{itemize}
\item Cl: $\bar m = 0.7577\times34.97 + 0.2423\times36.97 = 26.50 + 8.96 = \mathbf{35.45}$~amu
(closer to 35 because $\iso{35}{Cl}$ is more abundant).
\item Ar: almost pure $\iso{40}{Ar}$: $\bar m \approx 39.96$~amu (the 0.4\,\% of $\iso{36}{Ar}$ and $\iso{38}{Ar}$
lowers it to 39.95).
\item K: $\bar m = 0.9326\times38.96 + 0.0673\times40.96 = 36.33 + 2.76 = \mathbf{39.09}$~amu.
\end{itemize}
Conclusion: argon ($Z = 18$) is heavier than potassium ($Z = 19$). Ordering by mass, as Mendeleev did,
would put K before Ar, in the column of the noble gases, which is chemically absurd. The periodic table is
ordered by $Z$, not by mass.
\end{handout}

\sect{5.2}{Atomic radius}
\begin{concept}[Atomic radius.]
The distance from the nucleus to the outer edge of the electron cloud: the radius of the highest
occupied AO, i.e.\ the position of the maximum of its radial probability density. Since the density
fades out gradually, it is measured indirectly from distances between nuclei in molecules and solids.
\end{concept}
\begin{center}\small
\begin{tabular}{@{}l p{112mm}@{}}
\toprule
Covalent radius & half the distance between two identical bonded atoms ($\ce{H2}$, $\ce{Cl2}$ \dots)\\
Metallic radius & half the distance between neighbouring atoms in a metal crystal (Cu, Au)\\
Van der Waals radius & half the distance between two identical non-bonded atoms touching in a solid;
much larger than the covalent radius (used for noble gases, marked ``vdw'' in Table 5-1)\\
Ionic radius & from distances between ions in an ionic crystal; cation $<$ atom $<$ anion\\ \bottomrule
\end{tabular}
\end{center}
\begin{handout}{bullets and box of 5.2.2, trend of the atomic radius}
\begin{itemize}
  \item The radius \textbf{increases down a group} (Li 157, Na 191, K 235, Rb 250, Cs 272~pm).
  \item The radius \textbf{decreases across a period} (Li 157 $\to$ F 64~pm; Na 191 $\to$ Cl 99~pm).
        The largest atoms are at the bottom left (Cs), the smallest at the top right.
\end{itemize}
\begin{tabular}{p{0.46\linewidth}|p{0.46\linewidth}}
\textbf{Down a group:} each period adds a shell; the valence electrons are in an orbital of larger $n$
($r\propto n^2/Z_\mathrm{eff}$), while $Z_\mathrm{eff}$ barely changes. The atom gets bigger. &
\textbf{Across a period:} same valence shell $n$, but $Z_\mathrm{eff}$ increases (poor screening within
a shell), so the cloud is pulled closer to the nucleus. The atom gets smaller.
\end{tabular}\\[3pt]
The noble gas values (``vdw'') are van der Waals radii and cannot be compared with the others: that is
why Ne (154~pm) looks larger than F (64~pm).
\end{handout}
\begin{concept}[Ions.]
A cation is smaller than its atom (fewer electrons, less repulsion; often a whole shell lost:
$\ce{Na+}$ is $[\mathrm{Ne}]$). An anion is larger (extra electron, more repulsion, same nuclear charge).
In an isoelectronic series the radius decreases as $Z$ increases:
$r(\ce{O^2-}) > r(\ce{F-}) > r(\ce{Na+}) > r(\ce{Mg^2+})$ (all 10 electrons, 8 to 12 protons).
\end{concept}

\sect{5.3}{Ionisation energy}
\begin{concept}[(First) ionisation energy $E_I$.]
The minimum energy needed to remove the outermost electron from an atom in the gas phase:
$\mathrm X(\mathrm g)\to\mathrm X^+(\mathrm g) + \mathrm e^-$. Given in $\mathrm{kJ\,mol^{-1}}$ in chemistry
($1\ \mathrm{eV\ per\ atom} = 96.5\ \mathrm{kJ\,mol^{-1}}$). It is always positive. The second ionisation
energy removes a second electron from $\mathrm X^+$, and so on; each is larger than the previous one.
\end{concept}
\textbf{Measurement}: by atomic spectroscopy, from the series limit (for H, the Lyman limit at 91.2~nm gives
13.6~eV, Unit 2), or by photoelectron spectroscopy (PES): a photon of known energy $h\nu$ ejects an
electron whose kinetic energy $E_k$ is measured, and $E_I = h\nu - E_k$.
\begin{center}
\begin{tikzpicture}
\begin{axis}[width=0.92\linewidth, height=5.2cm, xlabel={$Z$}, ylabel={$E_I$ (kJ\,mol$^{-1}$)},
  xmin=0.5, xmax=20.5, ymin=0, ymax=2600, xtick={1,...,20}, tick label style={font=\scriptsize},
  label style={font=\small}, grid=major, grid style={black!10}]
\addplot[heading, thick, mark=*, mark size=1.4pt] coordinates {(1,1312)(2,2372)(3,520)(4,900)(5,801)
  (6,1086)(7,1402)(8,1314)(9,1681)(10,2081)(11,496)(12,738)(13,578)(14,787)(15,1012)(16,1000)(17,1251)
  (18,1521)(19,419)(20,590)};
\node[font=\tiny, above] at (axis cs:1,1312) {H};
\node[font=\tiny, above] at (axis cs:2,2372) {He};
\node[font=\tiny, above] at (axis cs:3,520) {Li};
\node[font=\tiny, above] at (axis cs:4,900) {Be};
\node[font=\tiny, above] at (axis cs:5,801) {B};
\node[font=\tiny, above] at (axis cs:7,1402) {N};
\node[font=\tiny, above] at (axis cs:8,1314) {O};
\node[font=\tiny, above] at (axis cs:10,2081) {Ne};
\node[font=\tiny, above] at (axis cs:11,496) {Na};
\node[font=\tiny, above] at (axis cs:12,738) {Mg};
\node[font=\tiny, above] at (axis cs:13,578) {Al};
\node[font=\tiny, above] at (axis cs:15,1012) {P};
\node[font=\tiny, above] at (axis cs:16,1000) {S};
\node[font=\tiny, above] at (axis cs:18,1521) {Ar};
\node[font=\tiny, above] at (axis cs:19,419) {K};
\end{axis}
\end{tikzpicture}\\
{\small First ionisation energies of the first 20 elements (Table 5-2; H: 1312).}
\end{center}
\begin{handout}{bullets and box of 5.3.2, trend of the ionisation energy}
Hydrogen, missing in Table 5-2: $E_I = 13.6$~eV $= 1312\ \mathrm{kJ\,mol^{-1}}$.
\begin{itemize}
  \item $E_I$ \textbf{increases across a period} (Li 520 $\to$ Ne 2081): $Z_\mathrm{eff}$ increases and
        the atom shrinks, so the electron is held more strongly. Noble gases have the highest values.
  \item $E_I$ \textbf{decreases down a group} (He 2372 $\to$ Rn 1037; Li 520 $\to$ Cs 376): the electron
        is further from the nucleus. Alkali metals have the lowest values.
  \item The increase across a period has two small \textbf{dips}: group 2 $\to$ 13 (Be 900 $>$ B 801,
        Mg 738 $>$ Al 578) and group 15 $\to$ 16 (N 1402 $>$ O 1314, P 1012 $>$ S 1000).
\end{itemize}
\begin{tabular}{p{0.46\linewidth}|p{0.46\linewidth}}
\textbf{Be $>$ B:} the electron removed from B is its single 2p electron, which is higher in energy (and
screened by 2s) than the 2s electron removed from Be. Be \obox{d}\;\obox{e,e,e} vs B \obox{d}\;\obox{u,e,e}. &
\textbf{N $>$ O:} N has a half-filled 2p (\obox{u,u,u}), each electron alone in its orbital. In O
(\obox{d,u,u}) one orbital holds a pair; their mutual repulsion makes one of them easier to remove.
\end{tabular}
\end{handout}
\begin{concept}[Successive ionisation energies.]
They increase slowly while valence electrons are removed, then jump enormously when the first core
electron is reached. The position of the jump gives the number of valence electrons. Mg: 738, 1451,
\textbf{7733}, 10\,543~$\mathrm{kJ\,mol^{-1}}$: the jump after the 2nd electron shows 2 valence electrons
(group 2).
\end{concept}

\sect{5.4}{Polarisability}
\begin{concept}[Polar entity and dipole moment.]
An entity is polar when the barycentres of its positive and negative charges do not coincide. It is
modelled by two charges $+q$ and $-q$ at a distance $d$, with dipole moment
$\boxed{\vec\mu = q\,\vec d}$, from $-q$ to $+q$ (the chemists' convention). SI unit: C\,m; in chemistry the
debye, $1\ \mathrm D = 3.336\times10^{-30}$~C\,m.
\end{concept}
\begin{concept}[Polarisability $\alpha$.]
How easily the electron cloud of an atom is distorted by an external electric field $\vec E$, creating an
induced dipole $\vec\mu_\mathrm{ind} = \alpha\vec E$. A large, diffuse, loosely held cloud is very
polarisable.
\end{concept}
\textbf{Measurement}: from the change of the refractive index of the substance in an electric field
(light itself is an oscillating field that distorts the clouds).
\begin{handout}{bullets of 5.4.2, trend of the polarisability}
\begin{itemize}
  \item Polarisability \textbf{increases down a group}: larger atoms, valence electrons further from the
        nucleus and better screened (I $>$ Br $>$ Cl $>$ F; Xe $>$ Kr $>$ Ar).
  \item It \textbf{decreases across a period}: the atoms shrink and $Z_\mathrm{eff}$ holds the cloud
        more tightly. It follows the atomic radius. Anions are more polarisable than cations.
\end{itemize}
\end{handout}
Consequence: large polarisable atoms attract each other more (London forces), which is why the boiling
points rise from $\ce{F2}$ (gas) to $\ce{I2}$ (solid) and from He to Rn.

\sect{5.5}{Electronegativity}
\begin{handout}{box of 5.5, bond A--A versus bond A--B}
\begin{tabular}{p{0.46\linewidth}|p{0.46\linewidth}}
\centering\textbf{A--A} (e.g.\ H--H, Cl--Cl) & \centering\textbf{A--B} (e.g.\ H--Cl)\tabularnewline
\centering
\begin{tikzpicture}
  \shade[ball color=heading!25, opacity=0.8] (0,0) ellipse (1.5 and 0.55);
  \fill (-0.6,0) circle (2pt) node[below=2pt, font=\scriptsize] {A};
  \fill (0.6,0) circle (2pt) node[below=2pt, font=\scriptsize] {A};
\end{tikzpicture} &
\centering
\begin{tikzpicture}
  \shade[ball color=heading!25, opacity=0.8] (0.25,0) ellipse (1.55 and 0.6);
  \fill (-0.6,0) circle (2pt) node[below=2pt, font=\scriptsize] {A};
  \fill (0.6,0) circle (2pt) node[below=2pt, font=\scriptsize] {B};
  \node[font=\scriptsize] at (-0.6,0.35) {$\delta+$}; \node[font=\scriptsize] at (0.75,0.35) {$\delta-$};
\end{tikzpicture}\tabularnewline
The shared electrons are equally attracted by both atoms: symmetric density, \textbf{non-polar} bond. &
B attracts the bonding electrons more than A: the density is shifted towards B, which carries a partial
charge $\delta-$ (A: $\delta+$). The bond is \textbf{polar} and has a dipole moment.
\end{tabular}
\end{handout}
\begin{concept}[Electronegativity $\chi$.]
The ability of an atom, \emph{in a bond}, to attract the shared electron density towards itself. It is a
relative number without unit. The larger $\chi_B - \chi_A$, the more polar the bond A--B; for a very large
difference ($\gtrsim 2$) the bond becomes ionic (NaCl).
\end{concept}
\subsect{5.5.1\quad Measurement: two scales}
\textbf{Pauling (experimental).} A polar bond A--B is \emph{stronger} than expected from the bonds
A--A and B--B, because the partial charges attract. The excess energy
$\Delta = D_{AB} - \sqrt{D_{AA}\,D_{BB}}$ (bond dissociation energies in $\mathrm{kJ\,mol^{-1}}$) gives
\[\boxed{|\chi_A - \chi_B| = 0.102\,\sqrt{\Delta}}\qquad\text{and the scale is anchored by }
\chi_\mathrm H = 2.20\ \text{(or } \chi_\mathrm F = 3.98).\]
\begin{worked}[: H--F from Table 5-3]
$\sqrt{436\times158} = 262.5$;\ $\Delta = 567 - 262.5 = 304.5\ \mathrm{kJ\,mol^{-1}}$;\
$|\chi_\mathrm F - \chi_\mathrm H| = 0.102\sqrt{304.5} = 1.78$, so $\chi_\mathrm F = 2.20 + 1.78 = 3.98$, the
value of Table 5-4. For H--I the same formula gives 0.66 (so $\chi_\mathrm I \approx 2.86$ instead of 2.66):
the tabulated values are a best fit over many molecules, and a single pair only gives an estimate.
(Pauling's original version used the arithmetic mean $\frac12(D_{AA} + D_{BB})$ and energies in eV:
$|\chi_A - \chi_B| = \sqrt{\Delta/(96.5\ \mathrm{kJ\,mol^{-1}})}$; use the version given in class.)
\end{worked}
\textbf{Allred--Rochow (theoretical).} Based on the electrostatic force felt by a valence electron at
the covalent radius $r$: $\chi = 0.359\,Z_\mathrm{eff}/r^2 + 0.744$ ($r$ in \AA). Large $Z_\mathrm{eff}$
and small radius mean high electronegativity. Both scales give very similar values.
\begin{handout}{bullets of 5.5.2, trend of the electronegativity}
\begin{itemize}
  \item $\chi$ \textbf{increases across a period} (Li 0.98 $\to$ F 3.98): $Z_\mathrm{eff}$ increases
        and the radius decreases. \textbf{Fluorine} is the most electronegative element, then O (3.44),
        then N and Cl.
  \item $\chi$ \textbf{decreases down a group} (F 3.98 $\to$ I 2.66; Li 0.98 $\to$ Cs 0.79): larger
        atoms hold bonding electrons less. Noble gases have no value (they hardly form bonds).
\end{itemize}
Metals (bottom left) have low $\chi$ and form cations; non-metals (top right) have high $\chi$.
\end{handout}

\begin{center}\small
\begin{tabular}{@{}l c c@{}}
\toprule
\textbf{Summary} & across a period $\rightarrow$ ($Z_\mathrm{eff}\nearrow$, same $n$) &
down a group $\downarrow$ ($n\nearrow$)\\ \midrule
atomic radius & decreases & increases\\
ionisation energy & increases & decreases\\
polarisability & decreases & increases\\
electronegativity & increases & decreases\\
atomic mass & increases & increases\\ \bottomrule
\end{tabular}
\end{center}
\begin{trap}
Do not mix up electronegativity (an atom \emph{in a bond} attracting shared electrons) and ionisation
energy (an \emph{isolated} atom losing an electron). Compare radii only within the same type
(covalent with covalent). ``Across a period'' trends apply to the s and p blocks; in the d block the
changes are much smaller.
\end{trap}

% ------------------------------------------------------------------ recap
\clearpage
\begin{recap}{Unit 5: atomic properties and periodicity}
\recapsub{Definitions}
\textbf{Atomic mass}: mass of one atom in amu; tabulated value = abundance-weighted average of the isotopes,
$\bar m = \sum x_i m_i$.\\
\textbf{Atomic radius}: radius of the highest occupied AO; measured as covalent, metallic, van der Waals
or ionic radius.\\
\textbf{Ionisation energy}: minimum energy for $\mathrm X(\mathrm g)\to\mathrm X^+(\mathrm g) + \mathrm e^-$,
in $\mathrm{kJ\,mol^{-1}}$.\\
\textbf{Polarisability}: ease of distortion of the electron cloud by an electric field; dipole
$\mu = qd$ (1~D $= 3.336\times10^{-30}$~C\,m).\\
\textbf{Electronegativity} $\chi$: ability of an atom in a bond to attract the shared electrons; no unit.

\recapsub{Trends}
\begin{tabular}{@{}l c c l@{}}
\toprule
 & across a period $\rightarrow$ & down a group $\downarrow$ & extremes\\ \midrule
radius & $\searrow$ & $\nearrow$ & Cs largest, He/F smallest\\
ionisation energy & $\nearrow$ (dips at 13 and 16) & $\searrow$ & He highest, Cs lowest\\
polarisability & $\searrow$ & $\nearrow$ & follows the radius\\
electronegativity & $\nearrow$ & $\searrow$ & F 3.98 highest, Cs/Fr lowest\\ \bottomrule
\end{tabular}\\[2pt]
Explanation: across a period $Z_\mathrm{eff}$ increases with the same $n$; down a group $n$ increases.

\recapsub{Formulas and numbers}
$\bar m = \sum x_i m_i$\quad(Cl 35.45, K 39.09, Ar 39.95).\quad
$1\ \mathrm{eV/atom} = 96.5\ \mathrm{kJ\,mol^{-1}}$.\quad PES: $E_I = h\nu - E_k$.\\
Pauling: $|\chi_A - \chi_B| = 0.102\sqrt{\Delta}$, $\Delta = D_{AB} - \sqrt{D_{AA}D_{BB}}$ (kJ\,mol$^{-1}$).\quad
Allred--Rochow: $\chi = 0.359\,Z_\mathrm{eff}/r^2 + 0.744$ ($r$ in \AA).

\recapsub{You should be able to}
\begin{itemize}
  \item compute an average atomic mass, or an abundance from an average mass;
  \item rank atoms and isoelectronic ions by radius, ionisation energy, polarisability or electronegativity,
        and justify with $n$ and $Z_\mathrm{eff}$;
  \item explain the Be/B and N/O anomalies with box diagrams;
  \item read successive ionisation energies to find the group;
  \item use the Pauling formula and judge the polarity of a bond.
\end{itemize}
\end{recap}
